\chapter{Conclusion}

This work consists of extending the solution of controlling a switched affine system 
subjected to dwell-time constraints to control this class of systems also with changing 
dynamics at switching times. This allows the method to be used in more control problems 
and maintains the convenience of solving the control problem with quadratic programming 
convex optimization, which is a mature class of optimization problems.

To test the method, first, a double integrator system with fixed dynamics was implemented in session \ref{cap4_aplication_integrador}
to validate the new solution. The result shows that the system diverges when not in the cyclic trajectory without the MPC controller and 
the controller is capable of correcting the trajectory and eventually engaging with the target trajectory.

The application example in session \ref{cap4_applicaton_patino1} introduces a system with changing dynamics at each switching time. 
A cyclic trajectory was calculated in session \ref{cap2_application_patino1} to use as an input parameter to implement the control method.
In this example, the circuit has stable dynamics that converge to the cyclic trajectory. Using the MPC controller reduces the time to converge
as can be seen in the comparison of the trajectories.

The multilevel converter in session \ref{cap4_applicaton_patino2} is a circuit that 
introduces more complexity, there are more operation modes and even modes that are repeated
during the same cycle. The simulation without the MPC shows a trajectory that converges to a cyclic limit 
that is different than the target cyclic trajectory. With the use of the MPC controller, we see that the system 
is driven to the target trajectory. The limit cycle of this circuit is also changing with different initial conditions,
in a first analysis this could be due to a numerical error in the simulation given that the numbers for this example 
are very small and susceptible to float point precision errors, but more work is needed to explain this behavior.

Future work will include the reimplementation of the linearization using augmented states, 
this approach was used by \citeonline{patino2010predictive} to write the system in an 
autonomous form and some initial analysis shows that this approach can simplify the 
linearization procedure to obtain the prediction equation for the cycle error and would 
contribute with improvement to the current work.

Another important work involves the physical laboratory implementation of a multilevel circuit to 
validate the control solution in an embedded environment; since the majority of existing 
literature remains in the simulation domain, a real-world application would significantly enrich 
the research. In line with this transition to practical implementation, future work contemplates 
incorporating discrete clock cycles into the simulation to account for the fact that switching 
instants must be multiples of the processor's clock period. For circuits with very small time 
constants, such as the multilevel converter, this discretization introduces rounding errors that 
can inject perturbations into the system, necessitating a detailed investigation into the 
resulting impact on stability and performance.

We propose the time to finish the project to be in three semesters. The first one is to 
finish the theoretical dependencies and analysis of the digital domain of the work, the 
second is to implement the real circuit in-lab and the last is to wrap all results in 
the final work.